\documentclass[11pt]{article}
\usepackage[a4paper,margin=25mm]{geometry}
\usepackage{amsmath,amssymb,amsthm,hyperref}
\newtheorem{theorem}{Theorem}
\newtheorem{lemma}{Lemma}
\newcommand{\PP}{\mathcal P}
\title{Refining pure densities onto one uniform rational grid}
\author{Complete proof; three independent mathematical reviews completed}
\date{30 September 2026}
\begin{document}\maketitle
Fix $q\ge0$. Here polynomial bounds are in a parameter $m$, and their
exponents may depend on $q$ and on the fixed input bounds. Rational
coefficients of the input densities need not have small denominators.

\begin{theorem}\label{main}
Suppose a partition of $[0,1]$ into rational intervals has polynomially
many cells, all endpoints on one polynomial rational grid, and all
cell lengths at least an inverse polynomial. On each cell $I$ there
is one positive rational polynomial density $f_I$ of polynomial degree,
with its maximum and reciprocal minimum bounded by polynomials.
Suppose all unnormalized cell moments through degree $q$ have one
polynomial common denominator.
For any polynomial degree bound $M\ge q$, there is a refinement of
this partition onto one uniform rational grid, and positive rational
polynomial densities on the refined cells, with these properties:
\begin{enumerate}
\item Every original cell moment through degree $M$ is unchanged.
\item Every refined cell moment through degree $q$ has one polynomial
common denominator, shared by all cells and all densities.
\item The new density is uniformly within an arbitrarily prescribed
inverse polynomial of the old density. In particular it is positive
and retains fixed conditional density-ratio bounds whenever the input
has such bounds.
\end{enumerate}
The uniform grid denominator, number of cells, and piece degrees are
polynomial. Labels of pure cells may be retained, so this refinement
preserves ordered degree-$M$ moments of any collection of densities
whose original cells are pure.
\end{theorem}

\begin{lemma}[A fixed-order polynomial quotient]\label{quotient}
Let $I\subset[0,1]$ be partitioned into intervals of common length
$h$. Let $P$ project onto $\PP_q$ separately on those intervals.
For $Q\in\PP_M$,
\[
 \operatorname{dist}_{L^2(I)}(Q,\PP_q)
 \le C_q(M+1)^{2q+2}h^{-q-1}\|(I-P)Q\|_{L^2(I)}.
\]
\end{lemma}
\begin{proof}
Subtract the degree-$q$ Taylor polynomial at the left endpoint and
integrate $Q^{(q+1)}$ repeatedly. Since $|I|\le1$, its $L^2$ norm
is bounded by $C_q\|Q^{(q+1)}\|_2$.
The $(q+1)$st derivative of $PQ$ is zero on every small interval.
The fixed-order polynomial inverse Markov estimate on each small
interval gives the stated bound after summing squares.
\end{proof}

\begin{proof}[Proof of Theorem~\ref{main}]
Choose one polynomial integer $Q$, divisible by the original endpoint
grid denominator and large enough that every original cell contains
at least $q+1$ intervals of length $h=1/Q$.
Choose a much finer polynomial integer $T$; its exponent will be fixed
last. On every small grid cell, round each of its raw moments through
$q$ to a multiple of $1/T$. Initially these differ from the original
moments by at most $1/T$.

We restore the first $q+1$ totals on each original cell $I$ without
introducing unrelated denominators. Use its first $q+1$ adjacent small
cells as controls. On these controls add constant-density corrections
$c_0,\ldots,c_q$ to the rounded moment vectors. Their moment matrix is
\[
 A_{rj}=\int_{a+jh}^{a+(j+1)h}x^r\,dx,
       \qquad 0\le r,j\le q,
\]
where $a$ is the left endpoint of $I$.
It is invertible: after $x=a+hy$, its determinant is
\[
 \det A=h^{(q+1)(q+2)/2}
       \det\left(\int_j^{j+1}y^r\,dy\right)_{r,j=0}^q.
\]
The latter determinant is a fixed nonzero rational number, because its
rows are monic polynomials of degrees $r$ evaluated at the distinct
integers $j$ (the leading coefficient of $\int_j^{j+1}y^rdy$ as a
polynomial in $j$ is one).
Thus $A^{-1}$ has norm bounded by a fixed power of $Q$, uniformly in
$a\in[0,1]$. More strongly, its entries have one denominator dividing
a fixed constant times a fixed power of $Q$: the determinant has the
displayed translation-independent form, and all entries of $A$ have
that property because $a$ is on the same grid.
The discrepancies of the coarse totals have one common denominator
dividing $T$ times the original low-moment denominator. Consequently
all corrected low moments, for every cell, have one polynomial common
denominator. Their errors from the original fine-cell moments are at
most $\operatorname{poly}(Q)/T$.

Fix an original cell $I$ and write $f$ for its original density.
Realize the small low-moment errors on each grid cell by a polynomial
$e$ of degree at most $q$. The rescaled Legendre moment basis gives
\[
 \|e\|_\infty\le C_q h^{-q-1}\operatorname{poly}(Q)/T.
\]
The correction has zero total moments through $q$ on $I$, by the
preceding restoration. We must restore its unwanted higher moments.
On $W=(I-P)\PP_M\subset L^2(I)$ define
\[
 \Lambda((I-P)F)=\int_I eF.
\]
This is well defined because $\ker((I-P)|_{\PP_M})=\PP_q$ and
$e$ annihilates $\PP_q$. Lemma~\ref{quotient} bounds its norm by
$C_q(M+1)^{2q+2}h^{-q-1}\|e\|_2$.
Let $z\in W$ be the Riesz representative. It has zero low moments
on every fine cell and
\[
 \|z\|_\infty
 \le (M+1)h^{-1/2}\|z\|_2
 \le \operatorname{poly}_q(M,Q)/T.
\]
The new density $f+e-z$ therefore has precisely the prescribed fine-cell
low moments and the original moments through $M$ on $I$.
Everything is rational: fine moments of $f$ are rational, the moment
systems are rational, and a rational basis of $W$ has rational Gram
matrix. Choosing the exponent and constant in $T$ sufficiently large
makes $\|e-z\|_\infty$ smaller than any prescribed inverse polynomial,
uniformly across all original cells. Positivity follows from the
assumed polynomial lower bound. The new polynomial degrees are at most
$\max(M,\deg f)$.

Every new-cell moment denominator is the already controlled target
denominator; the coefficients of $z$ may have much larger denominators,
which is harmless. If the original cells are pure for different labels,
a product of original cells fixes the relative order of all labels
that occupy distinct cells. At most one selected variable can occupy
each pure cell. Conditional order kernels of degree at most $M$ are
therefore preserved by the unchanged marginal cell moments. This proves
the final assertion.
\end{proof}
\end{document}
